\documentclass[12pt]{article}
\usepackage{amsmath}
\usepackage{physics}
\usepackage{geometry}
\geometry{a4paper, margin=1in}

\title{Detailed Quantum Mechanical Description of the Water ($H_2O$) Wave Function}
\author{}
\date{}

\begin{document}

\maketitle

\section{The Electronic Hamiltonian}
Under the Born-Oppenheimer approximation, the nuclear kinetic energy is neglected, and the nuclear-nuclear repulsion is treated as a constant. For a water molecule with 10 electrons and 3 nuclei (Oxygen at origin, Hydrogens at $\mathbf{R}_{H_A}$ and $\mathbf{R}_{H_B}$), the electronic Hamiltonian in atomic units is:

\begin{equation}
\hat{H}_{el} = -\sum_{i=1}^{10} \frac{1}{2} \nabla_i^2 - \sum_{i=1}^{10} \sum_{A=1}^{3} \frac{Z_A}{|\mathbf{r}_i - \mathbf{R}_A|} + \sum_{i=1}^{10} \sum_{j>i}^{10} \frac{1}{|\mathbf{r}_i - \mathbf{r}_j|}
\end{equation}

Where $Z_O = 8$ and $Z_{H_A} = Z_{H_B} = 1$. The terms represent electron kinetic energy, electron-nucleus Coulombic attraction, and electron-electron Coulombic repulsion, respectively.

\section{Basis Set and LCAO Approximation}
To solve the time-independent Schr\"{o}dinger equation $\hat{H}_{el} \Psi = E \Psi$, we expand the one-electron molecular orbitals (MOs), $\psi_i$, as a Linear Combination of Atomic Orbitals (LCAO):

\begin{equation}
\psi_i = \sum_{\mu=1}^K c_{\mu i} \phi_\mu
\end{equation}

Using a minimal basis set (e.g., STO-3G), we have $K=7$ basis functions:
\begin{itemize}
    \item Oxygen: $\phi_{O_{1s}}, \phi_{O_{2s}}, \phi_{O_{2p_x}}, \phi_{O_{2p_y}}, \phi_{O_{2p_z}}$
    \item Hydrogen A: $\phi_{H_A}$
    \item Hydrogen B: $\phi_{H_B}$
\end{itemize}

\section{Symmetry Adapted Linear Combinations (SALCs)}
Water belongs to the $C_{2v}$ point group. To block-diagonalize the Hamiltonian, we construct SALCs from the equivalent hydrogen $1s$ orbitals.
Assuming the molecule lies in the $yz$-plane with the $z$-axis bisecting the H-O-H angle:

\subsection{$a_1$ Symmetry (Totally Symmetric)}
These orbitals remain unchanged under all symmetry operations ($E, C_2, \sigma_v(xz), \sigma_v'(yz)$):
\begin{align}
\chi_1 &= \phi_{O_{1s}} \\
\chi_2 &= \phi_{O_{2s}} \\
\chi_3 &= \phi_{O_{2p_z}} \\
\chi_4 &= \frac{1}{\sqrt{2(1+S)}} (\phi_{H_A} + \phi_{H_B})
\end{align}
Where $S$ is the overlap integral $\langle \phi_{H_A} | \phi_{H_B} \rangle$.

\subsection{$b_2$ Symmetry}
Changes sign upon $C_2$ rotation and $\sigma_v(xz)$ reflection:
\begin{align}
\chi_5 &= \phi_{O_{2p_y}} \\
\chi_6 &= \frac{1}{\sqrt{2(1-S)}} (\phi_{H_A} - \phi_{H_B})
\end{align}

\subsection{$b_1$ Symmetry}
Changes sign upon $C_2$ rotation and $\sigma_v'(yz)$ reflection:
\begin{equation}
\chi_7 = \phi_{O_{2p_x}}
\end{equation}

\section{The Roothaan-Hall Equations}
To find the optimal coefficients $c_{\mu i}$, we minimize the total energy using the Variational Principle, leading to the pseudo-eigenvalue problem:

\begin{equation}
\mathbf{F} \mathbf{C} = \mathbf{S} \mathbf{C} \mathbf{\epsilon}
\end{equation}

Where:
\begin{itemize}
    \item $\mathbf{F}$ is the Fock matrix, with elements $F_{\mu\nu} = \langle \phi_\mu | \hat{f} | \phi_\nu \rangle$.
    \item $\mathbf{C}$ is the matrix of MO coefficients $c_{\mu i}$.
    \item $\mathbf{S}$ is the overlap matrix, $S_{\mu\nu} = \langle \phi_\mu | \phi_\nu \rangle$.
    \item $\mathbf{\epsilon}$ is a diagonal matrix of orbital energies.
\end{itemize}

Because $\hat{f}$ depends on the density matrix (and thus on $\mathbf{C}$), this equation is solved iteratively via the Self-Consistent Field (SCF) method.

\section{The Total Many-Body Wave Function}
For the closed-shell singlet ground state of $H_2O$ (${}^1A_1$), the 10 electrons occupy the 5 lowest-energy MOs ($1a_1, 2a_1, 1b_2, 3a_1, 1b_1$). The total N-electron wave function is written as a Slater determinant to enforce antisymmetry:

\begin{equation}
\Psi_{el} = \frac{1}{\sqrt{10!}} 
\begin{vmatrix}
\psi_{1a_1}(1)\alpha(1) & \psi_{1a_1}(1)\beta(1) & \dots & \psi_{1b_1}(1)\beta(1) \\
\psi_{1a_1}(2)\alpha(2) & \psi_{1a_1}(2)\beta(2) & \dots & \psi_{1b_1}(2)\beta(2) \\
\vdots & \vdots & \ddots & \vdots \\
\psi_{1a_1}(10)\alpha(10) & \psi_{1a_1}(10)\beta(10) & \dots & \psi_{1b_1}(10)\beta(10)
\end{vmatrix}
\end{equation}

\end{document}
